\documentclass[11pt]{article}

%------------------------- fonts & encoding -----------------------------------
\usepackage[T1]{fontenc}
\usepackage[utf8]{inputenc}
\usepackage{amsmath,amssymb,amsthm}
\usepackage{mathtools}
\usepackage[sc]{mathpazo}
\linespread{1.06}
\usepackage{microtype}

%------------------------- layout ---------------------------------------------
\usepackage[a4paper,margin=1.05in,marginparwidth=1.0in]{geometry}
\usepackage{booktabs}
\usepackage{enumitem}
\setlist{itemsep=2pt,topsep=3pt}
\usepackage{xcolor}
\usepackage{graphicx}
\usepackage{tikz}
\usepackage{pgfplots}
\pgfplotsset{compat=1.17}
\usepackage[most,breakable]{tcolorbox}

%------------------------- palette --------------------------------------------
\definecolor{crimson}{HTML}{A51C30}
\definecolor{ink}{HTML}{1A1A1A}
\definecolor{steel}{HTML}{2C5F8A}
\definecolor{amber}{HTML}{9A6A00}
\definecolor{teal}{HTML}{1F6F6B}
\definecolor{paperlt}{HTML}{F7F4EF}

%------------------------- section styling ------------------------------------
\usepackage{titlesec}
\titleformat{\section}{\Large\bfseries\color{steel}}{\thesection}{0.7em}{}
\titleformat{\subsection}{\large\bfseries\color{ink}}{\thesubsection}{0.6em}{}
\titlespacing*{\section}{0pt}{1.6em}{0.6em}

%------------------------- hyperref -------------------------------------------
\usepackage[colorlinks=true,linkcolor=steel,citecolor=steel,urlcolor=crimson]{hyperref}

%------------------------- theorem environments -------------------------------
\theoremstyle{definition}
\newtheorem{definition}{Definition}[section]
\newtheorem{example}{Example}[section]
\newtheorem{exercise}{Exercise}[section]
\theoremstyle{plain}
\newtheorem{proposition}{Proposition}[section]
\theoremstyle{remark}
\newtheorem*{remark}{Remark}

%------------------------- callout boxes --------------------------------------
\newtcolorbox{intuition}{
  enhanced, breakable, colback=crimson!5, colframe=crimson, boxrule=0.5pt,
  left=8pt,right=8pt,top=6pt,bottom=6pt, arc=2pt,
  colbacktitle=crimson, coltitle=paperlt, fonttitle=\bfseries, title={Intuition}}

\newtcolorbox{pitfall}{
  enhanced, breakable, colback=amber!7, colframe=amber, boxrule=0.5pt,
  left=8pt,right=8pt,top=6pt,bottom=6pt, arc=2pt,
  colbacktitle=amber, coltitle=paperlt, fonttitle=\bfseries, title={Pitfall}}

\newtcolorbox{keyresult}{
  enhanced, breakable, colback=steel!6, colframe=steel, boxrule=0.5pt,
  left=8pt,right=8pt,top=6pt,bottom=6pt, arc=2pt,
  colbacktitle=steel, coltitle=paperlt, fonttitle=\bfseries, title={Takeaway}}

\newtcolorbox{observed}{
  enhanced, breakable, colback=teal!6, colframe=teal, boxrule=0.5pt,
  left=8pt,right=8pt,top=6pt,bottom=6pt, arc=2pt,
  colbacktitle=teal, coltitle=paperlt, fonttitle=\bfseries, title={Measured}}

%------------------------- math operators -------------------------------------
\DeclareMathOperator{\Var}{Var}
\DeclareMathOperator{\Cov}{Cov}
\DeclareMathOperator{\Corr}{corr}
\DeclareMathOperator{\diag}{diag}
\DeclareMathOperator{\std}{std}
\DeclareMathOperator*{\avg}{\mathbb{E}}
\newcommand{\T}{^{\!\top}}
\newcommand{\R}{\mathbb{R}}
\newcommand{\half}{\tfrac12}
\newcommand{\wvec}{\mathbf{w}}
\newcommand{\Xmat}{\mathbf{X}}
\newcommand{\Fmat}{\mathbf{F}}
\newcommand{\Vmat}{\mathbf{V}}
\newcommand{\Umat}{\mathbf{U}}
\newcommand{\Rmat}{\mathbf{R}}
\newcommand{\Lam}{\boldsymbol{\Lambda}}
\newcommand{\Del}{\boldsymbol{\Delta}}
\newcommand{\fvec}{\mathbf{f}}

%------------------------- title ----------------------------------------------
\title{\vspace{-2.2em}
  {\color{steel}\rule{\textwidth}{1.3pt}}\\[0.7em]
  \Huge\bfseries Earnings Yield\\[0.35em]
  \large\mdseries The earnings-based value factor --- earnings per dollar of price, and a two-sided cross-section\\[0.4em]
  {\color{steel}\rule{\textwidth}{0.7pt}}}
\author{}
\date{}

%------------------------- watermark footer -----------------------------------
\usepackage{fancyhdr}
\pagestyle{fancy}
\fancyhf{}
\renewcommand{\headrulewidth}{0pt}
\renewcommand{\footrulewidth}{0pt}
\fancyfoot[L]{\footnotesize\color{black!45}MSCI Barra USE4 Imitation}
\fancyfoot[R]{\footnotesize\color{black!45}\thepage}

\begin{document}
\maketitle
\thispagestyle{empty}

\begin{abstract}
Earnings Yield (EYLD) is the earnings-based value factor: it measures earnings per dollar of price, the inverse of the price-earnings multiple, and so reads the value effect off the income statement rather than the balance sheet. USE4 defines EYLD as a blend of three earnings-to-price views -- a forward view from analysts' estimates (EPFWD), a cash-earnings view (CETOP), and a trailing view (ETOP) -- combined into a single descriptor and standardised across the estimation universe, with USE4 weighting the \emph{forward} view most heavily. Our personal build cannot source forward analyst estimates, so it drops that dominant forward view and renormalises over the two backward-looking views, leaning toward the cash-earnings measure; cash earnings themselves are proxied by adding depreciation and amortisation back to net income. Unlike most descriptors, EYLD is genuinely two-sided: a loss-making firm has negative net income and therefore a negative earnings yield, and the build keeps those names because a loss is real information. The result is a left-skewed cross-section with a tail of loss-makers -- precisely the population a risk model must not throw away.
\end{abstract}

\begin{tcolorbox}[colback=paperlt,colframe=ink!30,boxrule=0.4pt,arc=2pt,left=10pt,right=10pt,top=8pt,bottom=8pt]
\textbf{How to read these notes.} We move from \emph{why} Earnings Yield is a factor to \emph{how} USE4 specifies it, then to how we actually compute it from the data we have, with a deliberate stop at the one feature that makes EYLD unusual -- it can be negative. Read Section~\ref{sec:negative} slowly; it carries the weight of the set. The numbers labelled as measured come from a single run of our personal build.

\smallskip
Four coloured boxes recur: \textcolor{crimson}{\textbf{Intuition}} gives the mental picture, \textcolor{amber}{\textbf{Pitfall}} flags what breaks without the fix, \textcolor{steel}{\textbf{Takeaway}} states a load-bearing result, and \textcolor{teal}{\textbf{Measured}} reports real numbers from the build run.

\smallskip
\textbf{Reference material.}
\begin{itemize}[leftmargin=1.4em,topsep=2pt]
  \item External: Menchero, Orr \& Wang (2011), USE4 Empirical Notes (Appendix A) and Methodology Notes (Sec. 2.2--2.3).
\end{itemize}
\end{tcolorbox}

\tableofcontents
\bigskip

\section{What Earnings Yield is, and why it is a factor}\label{sec:what}

\paragraph{Learning objectives.}
\begin{itemize}[leftmargin=1.4em,topsep=2pt]
  \item State EYLD as earnings per dollar of price and relate it to the price-earnings multiple.
  \item Explain why earnings-to-price is a \emph{value} signal and how it differs from an asset-based value measure.
  \item Say in one sentence what the EYLD exposure is, before any standardisation.
\end{itemize}

Earnings Yield is value seen through the income statement. Where an asset-based value measure asks how much net worth you buy per dollar of price, EYLD asks how much \emph{earnings} you buy per dollar of price. It is the reciprocal of the familiar price-earnings multiple: a firm trading at twenty times earnings has an earnings yield of about five cents on the dollar, one trading at ten times earnings about ten cents. High earnings yield means the market is charging little for each dollar of current profit -- cheap by the earnings lens; low or negative earnings yield means the market is charging a great deal, or the firm is not currently profitable at all.

This is the same value effect that drives book-to-price, read off a different line of the financial statements. Asset-based value looks at accumulated net worth; earnings yield looks at the current flow of profit. The two usually agree on the cheap end of the market but can disagree sharply for firms whose book value and earnings tell different stories -- an asset-heavy firm earning little, or an asset-light firm earning a lot. USE4 carries both, and EYLD is the earnings half of that pair.

\begin{intuition}
The whole factor answers one question: \emph{how many cents of recent earnings am I buying for each dollar of price?} A high number means a lot of profit per dollar paid -- the textbook ``cheap'' stock. Hold that picture; every computational choice below is in service of measuring those cents cleanly.
\end{intuition}

\begin{keyresult}
The EYLD exposure is a standardised blend of earnings-to-price ratios. Everything that follows is about which earnings, which price, and how to standardise.
\end{keyresult}

\section*{Notation}
\addcontentsline{toc}{section}{Notation}

\begin{center}
\begin{tabular}{@{}ll@{}}
\toprule
Symbol & Meaning \\
\midrule
$E_i$ & trailing net income for stock $i$ (sum of four most recent quarterly filings) \\
$\text{D\&A}_i$ & trailing depreciation and amortisation for stock $i$ \\
$P_i$ & price (signal-date market capitalisation) for stock $i$ \\
$\text{ETOP}_i$ & trailing earnings-to-price, $E_i / P_i$ \\
$\text{CETOP}_i$ & cash earnings-to-price, $\approx (E_i + \text{D\&A}_i)/P_i$ \\
$\text{EPFWD}_i$ & forward earnings-to-price (USE4's heaviest-weighted view; dropped in this build) \\
$y_i$ & raw blended yield: a renormalised weighted average of $\text{CETOP}_i$ and $\text{ETOP}_i$, weighted toward cash \\
$x_i$ & standardised EYLD exposure for stock $i$ \\
$w_i$ & cap weight of stock $i$ in the ESTU \\
$\avg_{\text{cw}}, \std_{\text{ew}}$ & cap-weighted mean, equal-weighted standard deviation over the ESTU \\
$Q_1,\dots,Q_5$ & quintiles of raw $y_i$, from most negative ($Q_1$) to cheapest ($Q_5$) \\
ESTU & estimation universe (MSCI USA IMI proxy) \\
\bottomrule
\end{tabular}
\end{center}

\section{How USE4 defines it}\label{sec:def}

\paragraph{Learning objectives.}
\begin{itemize}[leftmargin=1.4em,topsep=2pt]
  \item Name the three earnings-to-price views USE4 blends, and which one it weights most.
  \item State the descriptor weight, the standardisation target, and the three-sigma trim.
\end{itemize}

USE4 builds Earnings Yield as a blend of three earnings-to-price descriptors (Menchero, Orr \& Wang 2011, Empirical Notes, App. A): a \emph{forward} view EPFWD (analysts' predicted earnings over price), a \emph{cash} view CETOP (cash earnings over price), and a \emph{trailing} view ETOP (realised trailing earnings over price). The three are not co-equal. USE4 weights the forward view most heavily -- it is the primary, dominant ingredient of the composite, carrying the bulk of the blend, with the cash and trailing views contributing the remainder. The three combine into a single Earnings Yield descriptor carried at descriptor weight $1.0$.

The descriptor is then standardised over the estimation universe to cap-weighted mean $0$ and equal-weighted standard deviation $1$ (Methodology Notes, Sec. 2.3), after a three-standard-deviation trim (Sec. 2.2). Fundamentals are aligned point-in-time by filing date, so the earnings used on any signal date are only those a market participant could actually have known. That is the entire specification: three views (one dominant), one descriptor, one standardisation.

\section{How we compute it: two holes, and a blend}\label{sec:compute}

\paragraph{Learning objectives.}
\begin{itemize}[leftmargin=1.4em,topsep=2pt]
  \item Explain why the forward view is dropped, that it was the heaviest-weighted view, and what that does to the factor.
  \item Reconstruct the cash-earnings proxy and the trailing-earnings sum.
  \item Identify the two computational choices that most change a stock's ranking.
\end{itemize}

The definition assumes three clean inputs. Our data supplies two of them, with a proxy, and not the third. Two holes, then a blend.

\paragraph{(a) No forward earnings.} The data carries no forward analyst estimates, so the EPFWD view cannot be computed and is dropped. This is the most material deviation in the whole factor, because the dropped view is not a minor third -- it is the one USE4 weights most heavily, the dominant component of the intended blend. The remaining two views are the minority ingredients of the original specification, and the build now rests entirely on them: it is renormalised over CETOP and ETOP, weighted toward the cash-earnings measure. This is a documented data limitation, not a modelling choice; the consequence is that our EYLD is \emph{purely backward-looking}. We record it plainly so that nothing downstream silently assumes a forward component that is not there.

There is a second, subtler consequence worth naming now, because it connects directly to Section~\ref{sec:negative}. The forward view is precisely the descriptor that could lift an expected-to-recover loss-maker back toward a positive yield: analysts' forward earnings can be positive even when trailing earnings are negative. Dropping it removes the one mechanism that would have pulled recovering loss-makers out of the negative left tail. So the backward-only build does not merely \emph{shrink} the factor -- it \emph{sharpens} the negative tail, making our EYLD more two-sided than full USE4 would be.

\paragraph{(b) Cash-earnings proxy.} There is no ready cash-earnings field either. We proxy it by adding back the single largest non-cash charge -- depreciation and amortisation -- to net income:
\[
\text{CETOP}_i \;\approx\; \frac{E_i + \text{D\&A}_i}{P_i},
\qquad
\text{ETOP}_i \;=\; \frac{E_i}{P_i}.
\]
The add-back is what separates the cash view from the trailing view; without it, the two would collapse into the same number.

\paragraph{(c) Trailing earnings, point-in-time.} The trailing earnings $E_i$ are summed from the four most recent quarterly filings whose datekey falls on or before the signal date -- one year of realised earnings -- and divided by the signal-date market capitalisation. The denominator $P_i$ is deliberately the \emph{current} (signal-date) market cap, not a filing-date or trailing-average price: this reprices stale earnings against today's valuation, which is exactly what makes E/P a live valuation signal even with backward-looking earnings. To tolerate slow filers we allow a staleness window of roughly a year and a half before a stock's earnings are treated as missing.

\begin{pitfall}
Two choices here move rankings more than they look. Pairing a \emph{stale} earnings number with the current price (or current earnings with a stale price) distorts $E/P$ for exactly the names whose price has just moved -- the failure mode the signal-date denominator makes possible and that you must watch. And omitting the cash add-back collapses CETOP into ETOP, discarding the part of the blend meant to be steadier. The quietest trap is that nothing in the output flags the missing forward view: the renormalisation silently rests the whole factor on backward data -- correct, but invisible unless stated out loud.
\end{pitfall}

\section{The two-sided signature: negative earnings are real}\label{sec:negative}

\paragraph{Learning objectives.}
\begin{itemize}[leftmargin=1.4em,topsep=2pt]
  \item Explain why EYLD can be negative when book-to-price and dividend yield effectively cannot.
  \item Justify keeping loss-makers (Figure~\ref{fig:eyld_dist}) rather than flooring earnings at zero.
  \item Describe the shape of the raw EYLD cross-section and where loss-makers sit.
\end{itemize}

This is the section that distinguishes EYLD from every other value descriptor. Book value is almost always positive; dividend yield is never negative -- a firm either pays or it does not. Earnings have no such floor. A loss-making firm has negative net income, and dividing a negative number by a positive price gives a \emph{negative earnings yield}. The factor is therefore genuinely two-sided, with a real left tail.

The build keeps these names. A loss is information, not an error to be cleaned away. A negative $E/P$ says something precise: the market is paying a positive price for a firm that is currently losing money, which is a bet on recovery -- exactly the kind of name whose risk profile a model needs to track. The convenience of a one-sided, all-positive signal would cost you the most informative part of the distribution. Flooring earnings at zero, or dropping loss-makers, would erase that left tail and hide the very population most likely to surprise.

\begin{intuition}
Picture two firms priced identically. One has $E/P \approx +0.06$ -- a going concern bought at a discount to its profits. The other has $E/P \approx -0.10$ -- a story stock whose price is untethered from current earnings, a bet on a future it has not yet earned. These are categorically different bets, not merely ``more cheap'' versus ``less cheap.'' The sign carries the distinction, which is why the build keeps it.
\end{intuition}

\begin{observed}
On the build run the raw EYLD cross-section is two-sided, with a heavy left tail of loss-makers; the post-trim ESTU distribution is \emph{left-skewed} (skew $\approx -1.4$). The ESTU median raw yield is $\approx 0.058$. Sorting into quintiles gives a strictly monotone profile from a \emph{negative} bottom quintile up to a cheap-by-earnings top: $Q_1 \approx -0.10$, $Q_2 \approx 0.03$, $Q_3 \approx 0.057$, $Q_4 \approx 0.084$, $Q_5 \approx 0.16$.
\end{observed}

\begin{figure}[htbp]
\centering
\includegraphics[width=\linewidth]{figures/eyld_dist.png}
\caption{The earnings-yield cross-section, pooled over in-ESTU stock-dates. Earnings yield is genuinely two-sided: about $16\%$ of names have \emph{negative} earnings yield (loss-makers, the shaded region left of the break-even line), giving a left-skewed distribution with a median near $0.06$. The build retains these negatives as real information rather than flooring them at zero.}
\label{fig:eyld_dist}
\end{figure}



\begin{pitfall}
Flooring negative earnings at zero, or dropping loss-makers entirely, discards the loss-maker tail -- precisely the names a risk model exists to watch. The fix for a noisy tail is a better transform, not deletion.
\end{pitfall}

\begin{intuition}
$E/P$ is not a clean ``cheapness'' dial near zero earnings. The instability is not in the level so much as in the \emph{ordering}: a firm a hair above breakeven shows an enormous yield, the same firm a hair below shows a large \emph{negative} yield, so a small accounting difference -- a one-time charge, a tax-timing item -- can flip a stock from one end of the cross-section to the other. The ratio is least trustworthy exactly where earnings cross zero. That is why USE4 blends the cash and (normally) forward views in the first place: they steady the signal where the trailing ratio is shakiest.
\end{intuition}

\begin{example}
Two firms each trade at a market cap of \$1,000m. Firm A earned \$50m last year; firm B lost \$30m. Then $\text{ETOP}_A = 0.05$ and $\text{ETOP}_B = -0.03$. If both carry \$40m of D\&A, their cash views are $\text{CETOP}_A = 0.09$ and $\text{CETOP}_B = 0.01$. The add-back lifts B from a reported loss to a small positive cash yield \emph{in this case} -- but only because D\&A exceeds the loss; a deeper loss, larger than D\&A, would leave even the cash view negative. The trailing view, meanwhile, still records B's reported loss. That contrast is exactly the nuance the cash view is meant to capture.
\end{example}

\section{Trim and standardisation}\label{sec:std}

\paragraph{Learning objectives.}
\begin{itemize}[leftmargin=1.4em,topsep=2pt]
  \item Apply the three-sigma trim and explain why both tails can bind for EYLD.
  \item State the standardisation target and read the stability diagnostics.
\end{itemize}

The raw blended yield $y_i$ is trimmed at three standard deviations. Because EYLD is two-sided, \emph{both} tails can bind -- the extreme cheap names on the right and the deep loss-makers on the left -- unlike a one-sided descriptor where only one tail clips. The trimmed values are then standardised over the ESTU to cap-weighted mean $0$ and equal-weighted standard deviation $1$:
\[
x_i \;=\; \frac{y_i - \avg_{\text{cw}}(y)}{\std_{\text{ew}}(y)},
\qquad
\avg_{\text{cw}}(y) = \frac{\sum_i w_i\, y_i}{\sum_i w_i}.
\]
The cap-weighted centring means the exposure is measured relative to the dollar-weighted market, while the equal-weighted scaling fixes the spread.

\begin{observed}
On the build run the trim clips $\approx 1.3\%$ of ESTU rows (both tails can bind). After standardisation the cap-weighted mean is $0$ to machine precision and the equal-weighted standard deviation is $1.0$. About $97\%$ of ESTU rows fall within $[\pm 3]$, but only $\approx 87.5\%$ of \emph{all} stocks do: the mean and standard deviation are fixed on the ESTU, so out-of-universe names -- exposed but not used to calibrate the standardisation -- can sit further out. Month-over-month rank stability is high, Spearman $\rho \approx 0.963$, though less sticky than a payout-based factor because earnings move more than payout policy. The quintile yields are monotone from $\approx -0.10$ ($Q_1$) to $\approx 0.16$ ($Q_5$), a spread of $\approx 0.26$. The deliverable is $\approx 1.54$M rows over 330 dates and $\approx 14.9$k tickers.
\end{observed}

\begin{exercise}
A firm reports trailing net income of \$20m, D\&A of \$60m, and a market cap of \$2,000m. Compute its ETOP and CETOP. Which view would flip sign if the firm had instead lost \$20m? (Hint: with these figures, only the trailing numerator can go negative, because D\&A here exceeds the loss; a larger loss would carry the cash view negative too.)
\end{exercise}

\begin{exercise}
Spearman $\rho \approx 0.963$ month-over-month for EYLD, noticeably below the near-perfect month-over-month stability typical of a payout-based factor. Explain in one or two sentences why earnings yield should rank-shuffle more than dividend yield. (Hint: think about which underlying quantity management can smooth and which it cannot.)
\end{exercise}

\begin{exercise}
Suppose you ``cleaned'' EYLD by flooring all negative earnings yields at zero before standardising. Sketch what happens to the cap-weighted mean, the left tail, and the bottom quintile. (Hint: re-read the observed $Q_1$ mean and ask where those firms go.)
\end{exercise}

\section{Summary, and one caveat}\label{sec:summary}

\paragraph{Learning objectives.}
\begin{itemize}[leftmargin=1.4em,topsep=2pt]
  \item Recite the EYLD build from filings to standardised exposure.
  \item State the two real limitations and the levers that would address them.
\end{itemize}

The build, end to end: sum trailing net income from the four most recent quarterly filings, point-in-time; add depreciation and amortisation back to form cash earnings; divide each by signal-date market cap to get the trailing and cash earnings-to-price views; blend the two, weighted toward cash earnings, since the forward view USE4 weights most heavily cannot be sourced; \emph{keep} negative earnings as real information; trim at three standard deviations; standardise to cap-weighted mean $0$, equal-weighted standard deviation $1$. The output is a two-sided, left-skewed, month-to-month stable value factor with a genuine loss-maker tail.

\begin{keyresult}
EYLD is earnings-based value with a real negative side. The two computations that matter most are the cash add-back (which separates CETOP from ETOP) and the decision to keep loss-makers (which preserves the left tail a risk model needs).
\end{keyresult}

\begin{pitfall}
Two honest limitations remain. \emph{First}, this EYLD is purely backward-looking -- the view USE4 weights most heavily has been dropped -- so it lags turnarounds and growth re-ratings, marking a recovering firm cheap or expensive a quarter or two late, and (as noted in Section~\ref{sec:compute}) leaving its loss-makers in the left tail that a forward view would have rescued. \emph{Second}, the $E/P$ ratio is ill-behaved near zero earnings: it is non-monotonic in ``cheapness,'' since a firm with tiny positive earnings shows an enormous yield and a tiny loss a large negative one, so the ordering of near-breakeven firms is noisy. The two levers are clear: a forward-earnings source to restore the dropped view, and a more robust earnings-yield transform that behaves smoothly through zero.
\end{pitfall}

\begin{remark}
None of this changes the factor's place in the model: EYLD remains the income-statement reading of value, sitting alongside an asset-based value measure. The caveats are about precision near the boundary and timeliness, not about whether the signal is real.
\end{remark}

\begin{center}\small\itshape
References: Menchero, Orr \& Wang (2011), \textit{The Barra US Equity Model (USE4)} --- Empirical Notes (Appendix A) and Methodology Notes (Sec. 2.2--2.3). ESTU treated as an MSCI USA IMI proxy; measured quantities are from a personal build run.
\end{center}

\end{document}
